% Editorial bootstrap from the immutable source revision.
% This file is canonical and may be edited independently.

\part{Analyticity and Obstructions to Totalized Elementary Expressions}


% Source fragment: Non-elementarity; prefix ne.
\section{Statement}

\subsection{The two inputs}

Write $\An(f)=\{x\in\RealNumbers : f \text{ is real analytic at } x\}$ for the
\emph{analytic locus} of $f\colon\RealNumbers\to\RealNumbers$.  The proof combines one negative
fact about the Fabius function with one positive fact about the explicitly
defined totalized elementary-expression class introduced below.  Throughout
this chapter, ``elementary'' is shorthand for that class; it is not a claim
about an unspecified larger differential-field convention.

\begin{theorem}[Nowhere analyticity; already formalized]\label{ne:thm:nowhere}
Let $F$ be the Fabius function, extended by $0$ to the left of $0$ and by $1$
to the right of $1$.  Then $\An(F)=\RealNumbers\setminus[0,1]$.  In particular $F$ is
real analytic at no point of $[0,1]$, endpoints included.
\end{theorem}

\begin{proof}
On $(0,1)$ this is the case $n=1$ of the forward-engine theorem
\Cref{ii:thm:main}; its complete derivative-spine proof is given in
\Cref{ii:sec:spines,ii:sec:proof-main}.  It remains only to account for the
totalization and the endpoints.  The function is constant, hence analytic,
on each component of $\RealNumbers\setminus[0,1]$.  If it were analytic at
$0$, its analytic germ would agree with the zero function on a left-hand
interval.  The identity theorem for real power series would then make it zero
on a two-sided neighbourhood of $0$, contrary to strict increase on
$[0,1]$.  At $1$ the same argument, applied to $1-F$, uses the constant
right-hand continuation.  Thus precisely the points outside $[0,1]$ are
analytic points.
\end{proof}

This is the unnumbered corollary to equation~(24) of Arias de Reyna,
\emph{Arithmetic of the Fabius function} (arXiv:1702.06487v3), transferred
from Rvachev's $\RvachevUp$ to $F$ and completed at the right endpoint.  The
same statement is machine checked as
\texttt{Fabius.canonical\_fabius\_analyticAt\_iff} in the module
\texttt{FabiusFunction.NowhereAnalytic}; the proof above records the human
argument and points to the full spine estimate used in its interior step.

\begin{theorem}[Density of the analytic locus]\label{ne:thm:density}
If $f\in\Elem$, then $\An(f)$ is a dense open subset of~$\RealNumbers$.
\end{theorem}

\begin{proof}
By \Cref{ne:lem:open}, the analytic locus of every real function is open, so
only density requires proof.  Induct on the derivation of $f$ in
\Cref{ne:def:elem}.  A constant and the identity are analytic everywhere.

Suppose first that the final constructor is binary.  By the induction
hypothesis, the analytic loci of its two arguments are dense and open; their
intersection is therefore dense.  On that intersection both arguments are
analytic, hence so are their sum and product.  Thus the analytic locus of the
result contains a dense set.  Negation is even more immediate: it preserves
every analytic point.

It remains to treat the unary constructors.  For the reciprocal, logarithm,
and a fixed real power, the outer function is analytic away from the single
value $0$; for the arcsine it is analytic away from $\{-1,1\}$; and for the
exponential, sine, cosine, and arctangent there are no exceptional values.
These assertions are \Cref{ne:lem:generators}.  In every case the exceptional
set is finite, so \Cref{ne:lem:key}, applied to the inner function supplied by
the induction hypothesis, makes the analytic locus of the composite dense.
These are all the constructors in \Cref{ne:def:elem}, and the induction is
complete.
\end{proof}

\begin{corollary}[Main result]\label{ne:cor:main}
Let $U\subseteq[0,1]$ have nonempty interior and let $g\in\Elem$.  Then $g$
and $F$ do not agree on $U$.  In particular no member of the totalized
expression class $\Elem$ agrees with $F$ on $(0,1)$.
\end{corollary}

\begin{proof}
By \Cref{ne:thm:density} the set $\An(g)$ is dense, so it meets the nonempty open
set $\operatorname{int}U$; pick $x\in\operatorname{int}U\cap\An(g)$.  If $g$
and $F$ agreed on $U$ they would have the same germ at $x$, because
$\operatorname{int}U$ is a neighbourhood of $x$ contained in $U$; analyticity
at a point depends only on the germ, so $F$ would be analytic at $x\in[0,1]$,
contradicting \Cref{ne:thm:nowhere}.
\end{proof}

\begin{remark}
Hypothesising nonempty \emph{interior} rather than openness costs nothing in
the proof and buys a little: $U$ may be a closed interval, or an interval with
a nowhere dense closed set removed, or any set that is merely fat somewhere.
(Removing \emph{all} the rationals would not do: the irrationals in an
interval have empty interior, and the hypothesis genuinely fails.)
What cannot be weakened is fatness itself: on a finite set an elementary
function certainly can agree with $F$, since a polynomial interpolates any
finite table of values.
\end{remark}

\begin{boxedremark}
Corollary~\ref{ne:cor:main} is the strengthened form used here of the assertion
``the Fabius function is not elementary'', and it is strictly stronger than the bare
non-elementarity of $F$ as a function on $\RealNumbers$.  The bare statement leaves
open the possibility that some elementary function agrees with $F$ on all of
$(0,1)$ and differs from it outside; Corollary~\ref{ne:cor:main} rules that out, and rules
it out on every subinterval.
\end{boxedremark}

\begin{remark}
Only a very weak consequence of \Cref{ne:thm:density} is actually used: that an
elementary function is analytic at \emph{at least one} point of every nonempty
open set.  Density is nevertheless the statement one must prove, because it is
the statement that survives the induction: ``analytic somewhere'' is not
preserved by intersection, whereas ``dense open'' is.
\end{remark}

\section{The totalized elementary-expression class}\label{ne:sec:elem}

\subsection{Definition}

For the theorem to have an exact scope, we fix a concrete grammar of total
real functions.  It contains the usual real expressions built from constants,
the identity, field operations, fixed powers, exponentials, logarithms, and
the listed trigonometric functions wherever those expressions are classically
defined.  It does not silently include arbitrary algebraic branches or
arbitrary interleaved algebraic--exponential--logarithmic towers; the precise
extensions proved later are stated separately.

\begin{definition}[Totalized elementary-expression class]\label{ne:def:elem}
The class $\Elem$ of \emph{totalized elementary expressions}
$\RealNumbers\to\RealNumbers$ is the smallest
class of functions such that
\begin{enumerate}
\item every constant function $x\mapsto c$ lies in $\Elem$, and the identity
      $x\mapsto x$ lies in $\Elem$;
\item if $f,g\in\Elem$ then $x\mapsto f(x)+g(x)$, $x\mapsto f(x)g(x)$ and
      $x\mapsto -f(x)$ lie in $\Elem$;
\item if $f\in\Elem$ then $x\mapsto f(x)^{-1}$ lies in $\Elem$;
\item if $f\in\Elem$ and $r\in\RealNumbers$ then $x\mapsto f(x)^r$ lies in $\Elem$;
\item if $f\in\Elem$ then each of
      $\exp\circ f$, $\log\circ f$, $\sin\circ f$, $\cos\circ f$,
      $\arcsin\circ f$, $\arctan\circ f$ lies in $\Elem$.
\end{enumerate}
\end{definition}

Two features of Definition~\ref{ne:def:elem} deserve comment.

\paragraph{Composition is a theorem, not an axiom.}
Definition~\ref{ne:def:elem} contains no clause ``$f,g\in\Elem\Rightarrow f\circ
g\in\Elem$''.  It does not need one.

\begin{proposition}\label{ne:prop:comp}
$\Elem$ is closed under composition.
\end{proposition}

\begin{proof}
Induct on the derivation of the outer function $f$.  Every clause of
Definition~\ref{ne:def:elem} commutes with precomposition by a fixed $g$: the constant
$x\mapsto c$ composed with $g$ is again that constant; the identity composed
with $g$ is $g$ itself, which is elementary by hypothesis; and
$(f_1+f_2)\circ g=(f_1\circ g)+(f_2\circ g)$, $(\exp\circ f)\circ
g=\exp\circ(f\circ g)$, and likewise for every remaining clause.
\end{proof}

Consequently $\Elem$ contains every function one expects: differences
$f-g=f+(-g)$, quotients $f/g=f\cdot g^{-1}$, natural powers, all polynomial
and rational functions, $\sqrt{\cdot}$ and the $n$-th roots,
the absolute value $|x|=\sqrt{x^2}$, the tangent $\sin/\cos$, the hyperbolic
functions $\sinh x=(\EulerE^x-\EulerE^{-x})/2$, $\cosh$, $\tanh$, the arccosine
$\pi/2-\arcsin$, and the inverse hyperbolic sine
$\operatorname{arsinh}x=\log\bigl(x+\sqrt{1+x^2}\bigr)$.

\paragraph{Total functions and junk values.}
Every member of $\Elem$ is a \emph{total} function $\RealNumbers\to\RealNumbers$.  This is a
deliberate choice, and it is the choice that makes Corollary~\ref{ne:cor:main} strongest.
The formalization inherits \texttt{Mathlib}'s conventions
\[
  0^{-1}=0,\qquad \log x=\log|x|,\qquad \log 0=0,\qquad
  0^{r}=0\ \ (r\ne 0),\qquad 0^{0}=1,
\]
\[
  x^{r}=|x|^{r}\cos(\pi r)\ \ (x<0),\qquad
  \arcsin x=\pm\tfrac{\pi}{2}\ \ (\pm x\ge 1),
\]
so that clauses (3)--(5) of Definition~\ref{ne:def:elem} never fail to produce a function.

\begin{boxedremark}
One might worry that junk values weaken the theorem, by admitting
``elementary'' functions that are not classically elementary.  They do the
opposite.  Admitting more functions into $\Elem$ makes the assertion ``$g\in
\Elem$ and $g=F$ on $U$'' easier to satisfy, hence its refutation
stronger.  And nothing is lost in the other direction: if a classical
elementary expression is well formed at every point of an open set $U$ --- no
vanishing denominator, no nonpositive logarithm, no real power of a
nonpositive base, no arcsine of a value outside $[-1,1]$ --- then on $U$ it
agrees pointwise with the total function produced by the same derivation,
since on $U$ the junk values are never consulted.  So every such classical
elementary function on $U$ is the restriction to $U$ of a member of $\Elem$,
and Corollary~\ref{ne:cor:main} applies to it.
\end{boxedremark}

\begin{warning}[Odd roots need a different derivation]\label{ne:warn:oddroot}
The condition ``no real power of a nonpositive base'' in the box above cannot
be dropped, and it is the one place where the total-function convention bites.
\texttt{Mathlib}'s $x^{r}$ at a negative base is $|x|^{r}\cos(\pi r)$, so
$(-8)^{1/3}=2\cos(\pi/3)=1$, whereas the classical real cube root of $-8$
is~$-2$.  The clause $x\mapsto f(x)^{r}$ of Definition~\ref{ne:def:elem} therefore supplies
the $n$-th roots only on $[0,\infty)$.

Nothing is lost, because the classical real odd root is in $\Elem$ by a
\emph{different} derivation,
\[
  x\ \longmapsto\ \frac{x}{|x|}\cdot|x|^{1/n},
\]
which uses only division, the absolute value and a real power, and which
evaluates to $0$ at $x=0$ under the same conventions.  This is
\texttt{IsElementary.signedRpow}, and
\texttt{Fabius.signedRoot\_pow} verifies formally that for odd $n$ its $n$-th
power is the identity, so it really is the real $n$-th root and not merely a
plausible substitute.
\end{warning}

\subsection{What is not claimed}\label{ne:sec:notclaimed}

$\Elem$ contains Mathlib's total real-power representatives of the principal
$n$-th roots on $[0,\infty)$ for positive $n$ and, by
Warning~\ref{ne:warn:oddroot}, the classical odd-root functions on all of $\RealNumbers$.
It is not, as defined, closed under passage
to an arbitrary algebraic function: Definition~\ref{ne:def:elem} has no clause
producing a continuous branch of $P(x,y)=0$ for a polynomial $P$ with
elementary coefficients whose Galois group is not solvable, and Liouville's
differential-field definition of ``elementary'' does admit such branches.
Section~\ref{ne:sec:algebraic} closes that gap, not by enlarging $\Elem$ but by
proving the conclusion of Corollary~\ref{ne:cor:main} directly for such branches.
What remains genuinely outside the result is stated there.

Two smaller gaps are worth naming here.

The first is that $\Elem$ is a class of \emph{total} functions, so a classical
elementary expression enters it only through a derivation whose junk values
are never consulted on the region of interest; Warning~\ref{ne:warn:oddroot}
records the one place where the obvious derivation is the wrong one.

The second concerns variable exponents.  Definition~\ref{ne:def:elem}
provides $x\mapsto f(x)^r$ for a \emph{fixed} real $r$, and
\texttt{IsElementary.rpow\_of\_pos} supplies $f^g$ whenever the base $f$ is
positive \emph{at every real} --- which covers $(1+x^2)^{\sin x}$ and any base
that is everywhere positive (no uniform lower bound is needed), but not
$x\mapsto x^x$, whose base is positive only on $(0,\infty)$.  That case is
handled instead by the region statement below.

The sign of an everywhere nonvanishing base is not an obstruction.
\texttt{Mathlib}'s $x^y$ is sign-dependent at a
negative base --- there $x^y=|x|^y\cos(\pi y)$ --- but a single elementary
formula covers both signs at once,
\[
  f^{g}\;=\;\exp(\log f\cdot g)\cdot
           \cos\!\Bigl(g\pi\cdot\tfrac{1-f/|f|}{2}\Bigr),
\]
since $(1-f/|f|)/2$ is the indicator of $f<0$, making the cosine $\cos 0=1$
where $f>0$ and $\cos(g\pi)$ where $f<0$; no absolute value is needed inside
the logarithm because \texttt{Mathlib} already sets $\log x=\log|x|$.  This is
\texttt{IsElementary.rpow\_of\_ne\_zero}, and it needs only that the base never
vanishes.  This is a sufficient uniform closure rule, not a claim that every
power whose base has a zero is non-elementary: fixed natural powers are
already covered by other constructors.  At
$f(x)=0$ the conventions $0^r=0$ for $r\ne0$ and $0^0=1$ are discontinuous in
the exponent, so the nonvanishing formula does not extend uniformly there.
Nothing is lost for Corollary~\ref{ne:cor:main}, and this is
formalized rather than asserted: \texttt{Fabius.eqOn\_rpow\_of\_pos}
states that on any set where the base is positive, $f^{g}$ agrees pointwise
with $\exp(\log f\cdot g)$, which is a member of $\Elem$ whenever $f$ and $g$
are.  Since Corollary~\ref{ne:cor:main} already rules out agreement on \emph{every}
nonempty open subset of $[0,1]$, and the base of a classical
variable-exponent expression is positive wherever the expression is well
formed, such an expression is excluded too.

We also make no claim about elementary functions of a complex variable, and
none about the closely related but distinct question of whether $F$ satisfies
an algebraic differential equation.

\section{Analyticity of the generators}

Each generator of Definition~\ref{ne:def:elem} is real analytic away from a finite set of
arguments.  All of these facts are in \texttt{Mathlib}, most of them in the
form ``$C^n$ for every $n$ in $\NaturalNumbers\cup\{\infty,\omega\}$''; specializing to the
analytic exponent $\omega$ and applying \texttt{ContDiffAt.analyticAt}
converts them to the statements below.

\begin{lemma}\label{ne:lem:generators}
\leavevmode
\begin{enumerate}
\item $\exp$, $\sin$, $\cos$ and $\arctan$ are real analytic at every point
      of $\RealNumbers$.
\item $\log$ is real analytic at every $x\ne 0$ --- at negative arguments as
      well, because $\log$ is even under the convention $\log x=\log|x|$.
\item For each fixed $r\in\RealNumbers$, the function $x\mapsto x^{r}$ is real analytic
      at every $x\ne 0$; at negative $x$ this uses the identity
      $x^{r}=|x|^{r}\cos(\pi r)$.
\item $\arcsin$ is real analytic at every $x\notin\{-1,1\}$; outside
      $[-1,1]$ it is locally constant.
\item $x\mapsto x^{-1}$ is real analytic at every $x\ne 0$.
\end{enumerate}
\end{lemma}

\begin{proof}
The four entire generators in (1) have convergent Taylor series at every real
centre.  On either component of $\RealNumbers\setminus\{0\}$ the totalized
logarithm is $\log|x|$, hence is analytic, proving (2).  For fixed $r$, on
$(0,\infty)$ the power is $\exp(r\log x)$, while on $(-\infty,0)$ it is the
constant multiple $\cos(\pi r)\exp(r\log(-x))$; this proves (3), including
the cases where the constant multiple vanishes.  On $(-1,1)$, $\arcsin$ is
analytic because its derivative $(1-x^2)^{-1/2}$ is analytic and it has the
usual convergent local primitive; on each exterior component the chosen
totalization is constant.  Only $\pm1$ remain exceptional, proving (4).
Finally $x\mapsto x^{-1}$ is represented near every nonzero centre by the
convergent geometric series
\[
 \frac1x=\frac1a\sum_{m\ge0}\left(-\frac{x-a}{a}\right)^m
 \qquad(|x-a|<|a|),
\]
which proves (5).
\end{proof}

\begin{remark}
The bad sets are sets of \emph{values} of the inner function, not sets of
points of the domain --- $\{0\}$ for the reciprocal, the logarithm and a real
power, $\{-1,1\}$ for the arcsine, and $\varnothing$ for the rest.  This is
the shape the induction consumes, and it is why Lemma~\ref{ne:lem:key} below is stated
with a finite set $B\subseteq\RealNumbers$ of forbidden values.
\end{remark}

\section{Density of the analytic locus}\label{ne:sec:density}

\subsection{Openness}

\begin{lemma}\label{ne:lem:open}
For every $f\colon\RealNumbers\to\RealNumbers$ the set $\An(f)$ is open.
\end{lemma}

\begin{proof}
Analyticity at $x$ means that $f$ agrees near $x$ with the sum of a
convergent power series; that series then converges on a whole ball around
$x$ and, after a change of origin, represents $f$ near every point of that
ball.  So $\An(f)$ contains a neighbourhood of each of its points.
\end{proof}

\subsection{The unary step}

The following lemma is where the argument is actually made.  Recall the shape
of the difficulty.  If $f$ is analytic on a dense open set and $g$ is analytic
everywhere except at finitely many values, one would like to say that
$g\circ f$ is analytic wherever $f$ is analytic and $f(x)$ avoids the bad
values.  That set need not be dense: consider $f\equiv 0$ and $g(y)=y^{-1}$.
What saves the statement is that in that degenerate case $g\circ f$ is
constant, hence analytic anyway.  Separating the two cases is exactly the
isolated-zeros theorem.

\begin{lemma}[Key lemma]\label{ne:lem:key}
Let $f,g\colon\RealNumbers\to\RealNumbers$ and let $B\subseteq\RealNumbers$ be finite.  Suppose $\An(f)$ is
dense and $g$ is analytic at every $y\notin B$.  Then $\An(g\circ f)$ is
dense.
\end{lemma}

\begin{proof}
Let $U$ be a nonempty open set; we must produce a point of $U\cap\An(g\circ
f)$.  Since $\An(f)$ is dense there is $x_0\in U\cap\An(f)$, and by
Lemma~\ref{ne:lem:open} the set $V=U\cap\An(f)$ is an open neighbourhood of $x_0$.
Distinguish two cases.

\emph{Case 1: $f$ is constant equal to some $c\in B$ on a neighbourhood of
$x_0$.}  Then $g\circ f$ equals the constant $g(c)$ on that neighbourhood, so
$g\circ f$ is analytic at $x_0$ itself, and $x_0\in U$.

\emph{Case 2: for no $c\in B$ is $f$ eventually equal to $c$ near $x_0$.}  Fix
$c\in B$.  The function $f-c$ is analytic at $x_0$, so by the isolated-zeros
theorem for analytic functions of one real variable, either $f-c$ vanishes
identically near $x_0$ --- excluded in this case --- or $f-c$ is nonvanishing
on some punctured neighbourhood of $x_0$.  Thus for each of the finitely many
$c\in B$ there is a punctured neighbourhood of $x_0$ on which $f\ne c$.
Intersecting these finitely many punctured neighbourhoods with the punctured
neighbourhood $V\setminus\{x_0\}$ leaves a punctured neighbourhood of $x_0$,
which is nonempty because $\RealNumbers$ has no isolated points.  Any point $t$ of it
satisfies $t\in U$, $f$ analytic at $t$, and $f(t)\notin B$; hence $g$ is
analytic at $f(t)$ and $g\circ f$ is analytic at $t$.
\end{proof}

\begin{remark}
The case distinction is not an artefact of the write-up.  Without it the lemma
is false as stated, for the naive reason above.  What makes the two cases
exhaustive is the isolated-zeros theorem, applied to $f-c$: an analytic germ
either vanishes identically or vanishes nowhere on some punctured
neighbourhood.  A merely $C^\infty$ inner function admits neither
alternative --- it can equal $c$ on a sequence accumulating at $x_0$ without
being constant there --- so analyticity of $f$, and not just smoothness, is
what the step consumes.  Within the induction of this section it is the only
place where a theorem of substance about real analytic functions is used;
\Cref{ne:sec:algebraic,ne:sec:inverse} call on the implicit function theorem as
well.
\end{remark}

\subsection{Assembly}

The complete structural induction proving \Cref{ne:thm:density} is placed
beside the theorem.  The present section supplies precisely its two
non-formal ingredients: \Cref{ne:lem:open} makes the invariant open, while
\Cref{ne:lem:key} shows that a unary constructor with finitely many exceptional
input values preserves density.  Together with the exceptional sets computed
in \Cref{ne:lem:generators}, this exhausts the grammar of
\Cref{ne:def:elem}; there is no second proof to add here.

\begin{remark}[Sharpness]
Density cannot be strengthened to ``$\An(f)=\RealNumbers$'': the elementary functions
$x\mapsto x^{-1}$, $x\mapsto|x|$ and $x\mapsto\arcsin x$ have analytic loci
$\RealNumbers\setminus\{0\}$, $\RealNumbers\setminus\{0\}$ and $\RealNumbers\setminus\{-1,1\}$.  Nor can
``dense open'' be strengthened to ``open with finite complement'': the
elementary function $x\mapsto (\sin x)^{-1}$ omits the infinite discrete set
$\pi\IntegerNumbers$.  What density does give for free is that the complement is closed
and nowhere dense, being the complement of a dense open set.  It need not be
\emph{discrete}, and one more reciprocal settles that: $x\mapsto(\sin(1/x))^{-1}$
is elementary --- it is \texttt{IsElementary.id.inv.sin.inv} --- and its
analytic locus omits every $1/(k\pi)$ with $k\in\IntegerNumbers\setminus\{0\}$, together
with $0$, where the function is not even continuous.  Those points accumulate
at $0$.  So ``nowhere dense'' is the correct description of the exceptional
set and cannot be improved to ``discrete''.
\end{remark}

\section[Algebraic branches]{Algebraic functions over the elementary functions}\label{ne:sec:algebraic}

Definition~\ref{ne:def:elem} reaches principal $n$-th roots on nonnegative inputs
for positive $n$, and classical odd roots on all of $\RealNumbers$, but not a general
algebraic function.  This section removes that restriction.  It does so without
enlarging $\Elem$, for a reason worth stating: an algebraic clause inside the
inductive definition would cost Proposition~\ref{ne:prop:comp}.  Composing a
branch $y$ with an elementary $h$ produces a branch over the coefficients
$a_i\circ h$, but it is a \emph{continuous} branch only if $h$ is continuous,
and elementary functions need not be --- $x\mapsto x^{-1}$ is not.  Rather than
weaken the closure theorem, we prove the conclusion directly.

\begin{theorem}[Branches of polynomial equations]\label{ne:thm:algebraic}
Let $U\subseteq\RealNumbers$ be open, let $a_0,\dots,a_n$ be real analytic on $U$ with
$a_n$ nowhere zero on $U$, and let $y$ be continuous on $U$ with
\[
  \sum_{i=0}^{n} a_i(x)\,y(x)^i \;=\; 0 \qquad (x\in U).
\]
Then $y$ is real analytic at some point of every nonempty open subset of $U$.
\end{theorem}

Density is again the right conclusion, and for the same kind of reason as
before: $y=|\cdot|$ is a continuous branch of $y^2-x^2=0$ on all of $\RealNumbers$, with
$a_0=-x^2$, $a_1=0$, $a_2=1$, and it is not analytic at the origin.

\begin{proof}
The implicit function theorem applies only where $\partial P/\partial y\ne 0$,
and the classical way to control the locus where it vanishes is the
discriminant.  \texttt{Mathlib} has no usable discriminant theory ---
\texttt{Polynomial.discr} carries no bridge to separability, and there is no
continuity of roots in the coefficients --- so we induct on the degree
instead, which needs neither.

Write $g$ for $\partial P/\partial y$ evaluated along the branch and split
$U$ into $V=\{x\in U: g(x)\ne0\}$, which is open, and
$W=U\setminus\overline V$, which is also open.  On $V$ the implicit function
theorem applies directly.  On $W$ we have $g\equiv 0$, so $y$ satisfies the
\emph{derivative} equation, of degree $n-1$, whose leading coefficient
$n\,a_n$ is still nowhere zero; the inductive hypothesis applies to it.
Finally $V$ together with a dense subset of $W$ is dense in $U$, because an
open set disjoint from $V$ is disjoint from $\overline V$ and hence contained
in $W$.  The set $W$ absorbs the \emph{interior} of the degenerate locus
$\{g=0\}$; what the argument gives up is the boundary $\partial V\cap U$, and
that is exactly why the conclusion is density rather than analyticity
everywhere.  For $y=|\cdot|$ the degenerate locus is $\{0\}$, so $V=\RealNumbers
\setminus\{0\}$, $\overline V=\RealNumbers$, $W=\varnothing$, and the origin is conceded
to the splitting step --- correctly, since $|\cdot|$ is not analytic there.

The induction terminates because the degree drops.  Its base case $n=0$ is
vacuous rather than positive: the equation reads $a_0(x)=0$ while the leading
coefficient hypothesis says $a_0(x)\ne0$, so the region is empty.

The implicit function step is \texttt{Mathlib}'s
\path{ContDiffAt.implicitFunction}, which is stated for every smoothness
exponent in $\NaturalNumbers\cup\{\infty,\omega\}$; instantiating at $\omega$ and converting
with \path{ContDiffAt.analyticAt} gives an analytic implicit function
theorem.  Its one nontrivial hypothesis, invertibility of the partial
derivative as a continuous linear map, is discharged in one variable by
\path{HasDerivAt.hasFDerivAt_equiv} together with uniqueness of the
Fréchet derivative.
\end{proof}

\begin{corollary}\label{ne:cor:algebraic}
Let $U\subseteq[0,1]$ be nonempty and open, let $a_0,\dots,a_n\in\Elem$,
let $a_n$ vanish nowhere \emph{on $U$}, and let $y$ be continuous \emph{on
$U$} with $\sum_{i\le n}a_i(x)y(x)^i=0$ for every $x\in U$.  Then $y$ does not
agree with the Fabius function on $U$; in particular not on $(0,1)$.
\end{corollary}

\begin{proof}
Each $a_i$ has dense open analytic locus by \Cref{ne:thm:density}.  The
intersection $W$ of the $n+1$ of them is a finite intersection of dense open
sets, hence dense and open, so $U\cap W$ is nonempty and open, and
\Cref{ne:thm:algebraic} applies to it: $y$ is analytic at some $x\in U\cap
W\subseteq U$.  Now argue as in Corollary~\ref{ne:cor:main}.
\end{proof}

\begin{boxedremark}
Every hypothesis on the branch in Corollary~\ref{ne:cor:algebraic} is confined to
$U$, and that is what makes the corollary say what one wants it to say.  A
version with global hypotheses would be much weaker than it looks.  For the
standard Bring equation $y^5-y-x=0$, a global continuous branch would be a
continuous right inverse of $h(w)=w^5-w$.  Such a right inverse is injective;
its image is an interval, and $h(y(x))=x$ forces that image to be unbounded in
both directions, hence all of $\RealNumbers$.  It would therefore force $h$ itself to be
injective, a contradiction.  The localized statement nevertheless reaches
the positive real branch over $(0,1)$, an algebraic function not expressible by
radicals.

The corollary is genuinely stronger than Corollary~\ref{ne:cor:main} in that
direction, and independent of it in the other: it says nothing about $\exp$ or
$\log$ of an algebraic function, because it is a single algebraic step over
elementary coefficients rather than a closure.  A tower that interleaves
algebraic extensions with exponentials and logarithms at several levels is
covered only when it can be presented in one of the two forms.
\end{boxedremark}

\section{The inverse, and why inverting cannot help}\label{ne:sec:inverse}

$F$ is a homeomorphism of $[0,1]$ onto itself, so it has a continuous inverse.
We write $F^{-1}\colon\RealNumbers\to\RealNumbers$ for the formal development's clamped
totalization: the inverse on $[0,1]$, the constant $0$ on $(-\infty,0]$, and
the constant $1$ on $[1,\infty)$.  Everything above concerned $F$; this
section treats that totalization, and then closes the class under the
admissible inverse-branch rule.

In the formalization this function is \texttt{Fabius.fabiusInv}, defined by
\texttt{Set.IccExtend} of the order isomorphism.  The extension is continuous,
so nothing is hidden in the gluing; its constant tails are what the endpoint
half of \Cref{ne:thm:invnowhere} exploits.

\subsection{A right inverse is an implicit branch}

The observation that organizes this section is that no new analysis is
needed.  The equation
\[
  h\bigl(g(x)\bigr)=x
\]
says exactly that $g$ is a continuous branch of $P(x,z)=h(z)-x=0$, an implicit
equation whose partial derivative in $z$ is $h'$.  So the analytic implicit
function theorem quoted in \Cref{ne:sec:algebraic} --- already proved there, for
an unrelated purpose --- delivers the analytic inverse function theorem with
nothing further to check.

\begin{theorem}[Analytic inverse function theorem]\label{ne:thm:ift}
If $h$ is analytic at $g(x_0)$ with $h'(g(x_0))\ne0$, $g$ is continuous at
$x_0$, and $h\circ g$ is the identity near $x_0$, then $g$ is analytic
at~$x_0$.
\end{theorem}

\begin{proof}
Set $H(x,z)=h(z)-x$.  It is analytic near $(x_0,g(x_0))$, satisfies
$H(x_0,g(x_0))=0$, and has
$\partial_zH(x_0,g(x_0))=h'(g(x_0))\ne0$.  The analytic implicit-function
theorem therefore gives an analytic function $\widetilde g$ near $x_0$ with
$h(\widetilde g(x))=x$ and $\widetilde g(x_0)=g(x_0)$.  After shrinking the
neighbourhood, $h$ is one-to-one near $g(x_0)$ by the ordinary inverse-function
theorem; continuity of $g$ keeps $g(x)$ in that same neighbourhood.  The two
right inverses must then agree, $g=\widetilde g$, so $g$ is analytic at $x_0$.
\end{proof}

\subsection{The inverse of the Fabius function}

\begin{theorem}\label{ne:thm:invnowhere}
$F^{-1}$ is real analytic at no point of $[0,1]$; its analytic locus is
exactly $\RealNumbers\setminus[0,1]$.
\end{theorem}

\begin{proof}
In the interior, run \Cref{ne:thm:ift} backwards.  If $F^{-1}$ is analytic at
$y_0\in(0,1)$ then $F$, being its inverse near that point, is analytic at
$F^{-1}(y_0)\in(0,1)$ --- provided $(F^{-1})'(y_0)\ne0$.  It has no critical
point anywhere in the interior.  This is the one place where $F'$ enters at
all --- everywhere else
$F$ is used through its continuity and through its failure to be analytic ---
and what is needed of it is only that $F'>0$ on $(0,1)$.  Given that, the easy
half of the inverse function theorem --- the half that assumes the inverse
already exists --- yields
 \[
   (F^{-1})'(y)=\frac{1}{F'\bigl(F^{-1}(y)\bigr)}\qquad(0<y<1),
 \]
 a nonzero real number.  The core inverse module formalizes this reciprocal
 derivative, its strict positivity, and the stronger fact that $F^{-1}$ is
 $C^\infty$ throughout $(0,1)$.  \Cref{ne:thm:nowhere} then supplies the
 contradiction.

The direction deserves to be stated, because the shorter route is circular.
Differentiating $F^{-1}\circ F=\mathrm{id}$ to obtain
$(F^{-1})'\bigl(F(x)\bigr)\cdot F'(x)=1$ presupposes that $F^{-1}$ is
differentiable, which is exactly what is at issue; strict positivity of $F'$
is what supplies it.

At the endpoints the argument is different, and simpler.  $F^{-1}$ is
identically $0$ to the left of $0$; an analytic germ agreeing with $0$ on a
set with an accumulation point vanishes identically, so $F^{-1}$ would vanish
just to the right of $0$ as well, contradicting strict monotonicity.
Symmetrically at $1$.
\end{proof}

\begin{corollary}\label{ne:cor:invmain}
No member of $\Elem$ agrees with $F^{-1}$ on any subset of $[0,1]$ with
nonempty interior; in particular the restriction of $F^{-1}$ to $(0,1)$ has
no representative in this totalized expression class.
\end{corollary}

\begin{proof}
If a member $g$ of $\Elem$ agreed with $F^{-1}$ on a set $U\subseteq[0,1]$
having nonempty interior, density in \Cref{ne:thm:density} would provide
$y\in\operatorname{int}U$ at which $g$ is analytic.  Equality on the
neighbourhood $\operatorname{int}U$ gives $g$ and $F^{-1}$ the same germ at
$y$, making $F^{-1}$ analytic there and contradicting
\Cref{ne:thm:invnowhere}.
\end{proof}

\begin{boxedremark}
One might have hoped that inverting improves matters.  $F$ is flat to infinite
order at the endpoints, so $F^{-1}$ is steep there, and steepness is not
obviously an obstruction to being elementary --- the signed real cube root
$x\mapsto (x/|x|)|x|^{1/3}$ is steep at the origin and perfectly elementary.
\Cref{ne:thm:invnowhere} says the hope
is misplaced, and the reason is structural rather than quantitative: analytic
functions are closed under inversion wherever the derivative survives, so an
analytic inverse would hand analyticity straight back.
\end{boxedremark}

\subsection{Closing the class under inversion}

Nothing about the argument was special to $F$.  The same reasoning shows that
the invariant of \Cref{ne:thm:density} survives the admissible inverse-branch rule
below.

\begin{theorem}[Dense analytic locus of a right inverse]\label{ne:thm:invbranch}
Let $f$ be analytic on a dense subset of $\RealNumbers$, and let $g$ be
continuous on a nonempty open set $U$ with $f(g(y))=y$ there.  Then
$\An(g)\cap U$ is relatively open and dense in $U$.
\end{theorem}

\begin{proof}
Relative openness follows from \Cref{ne:lem:open}; we prove density.  Let
$I\subseteq U$ be any nonempty open interval.  The right-inverse identity
makes $g$ injective on $I$.  A continuous injective real function on an
interval is strictly monotone, so $J=g(I)$ is a nonempty open interval and
$g:I\to J$ is a homeomorphism.  Moreover $f|_J$ is its inverse and is
therefore injective.

The open set $J$ meets the dense analytic locus of $f$.  Choose a nonempty
open interval $V\subseteq J$ on which $f$ is analytic.  Since $f$ is
injective on $J$, it is not constant on $V$.  Hence the analytic function
$f'$ is not identically zero on $V$; by the isolated-zero theorem, some
$z\in V$ satisfies $f'(z)\ne0$.  Put $y=f(z)\in I$.  Near $y$ we have
$f\circ g=\id$, and $g(y)=z$; \Cref{ne:thm:ift} shows that $g$ is analytic at
$y$.  Thus every nonempty open interval in $U$ meets $\An(g)$, as required.
\end{proof}

Applied to every nonempty open subset at once, this says that a continuous
inverse branch is analytic on a relatively dense subset of its domain, and the
domain need not be open:
\path{Fabius.analyticDenseOn_of_rightInverse}.  For a
domain with empty interior the statement is vacuous, which is the correct
reading --- there is no open set on which to test.

The conclusion is ``analytic somewhere in $U$'', not ``analytic on $U$'', and
that is not slack in the argument: the real cube root is a continuous right
inverse of the entire function $x\mapsto x^3$ on all of $\RealNumbers$ and is not
analytic at the origin.  In the formalization the right inverse has to be
written $t\mapsto(t/|t|)\,|t|^{1/3}$, since \texttt{Real.rpow} makes
$(-8)^{1/3}=1$ (Warning~\ref{ne:warn:oddroot}); that this signed form really does
cube back to $t$ is \texttt{Fabius.signedRoot\_pow}.  Critical points of $f$
do produce singularities of the inverse.  It is the same phenomenon that made
density, rather than analyticity everywhere, the right invariant in
\Cref{ne:thm:density}.

\begin{definition}\label{ne:def:elemplus}
$\Elem^{-1}$ is the smallest class containing $\Elem$, closed under the
operations of Definition~\ref{ne:def:elem}, and closed under the following rule: if
$f\in\Elem^{-1}$, $U$ is open, $g$ is continuous on $U$ with $f(g(y))=y$ for
$y\in U$, and $g$ is analytic on the interior of the complement of $U$, then
$g\in\Elem^{-1}$.
\end{definition}

Two features of the rule are deliberate, and the Lambert $W$ function shows
why each is needed.  It is local in $U$: $W$ satisfies its defining identity
$W(y)\EulerE^{W(y)}=y$ only on $[-1/\EulerE,\infty)$, so a rule demanding a global
identity would not reach it.  And it asks for analyticity on the
\emph{interior} of the complement rather than at every point outside $U$: the
branch point $-1/\EulerE$ sits on the boundary of the natural $U=(-1/\EulerE,\infty)$,
and $W$ is certainly not analytic there.  Dropping that boundary costs
nothing, since the boundary of an open set is nowhere dense, and for $W$ the
constant extension then satisfies the analyticity clause on all of
$(-\infty,-1/\EulerE)$.

The Lean result is a conditional criterion for any supplied total function
satisfying these hypotheses.  Since \texttt{Mathlib} does not define $W$, the
formalization neither constructs a standard real branch nor verifies a
particular totalization of it.

That clause is a condition on the chosen extension of the branch past $U$, not
a fact that comes for free.  Extension by a constant always meets it, and that
is the honest reading of the rule: a branch is admitted together with a
declaration of how it is continued.  Other conventions also meet it without
being constant --- \texttt{Mathlib} sets $\log y=\log|y|$, which is analytic
off the origin.  At the other extreme $U=\emptyset$ reduces the clause to
``$g$ is entire'', so the rule also adjoins every everywhere-analytic
function; harmlessly, since those already have dense analytic locus, which is
all the invariant asks.  Inversion may be applied at any depth: to members of
$\Elem$, to branches of those, and so on without limit.

One asymmetry with $\Elem$ deserves stating, since \Cref{ne:sec:elem} made a
point of it.  Composition is a \emph{theorem} for $\Elem$
(Proposition~\ref{ne:prop:comp}), proved by inducting on the derivation of the
outer function, because every clause of Definition~\ref{ne:def:elem} commutes with
precomposition.  That induction does not replay for $\Elem^{-1}$: the inverse
clause does not commute with precomposition, since $g$ being a branch of
$f(z)=x$ says nothing about $g\circ h$ being a branch of anything.  So
$\Elem^{-1}$ is not claimed to be closed under composition, and
\Cref{ne:thm:plusdense} does not need it --- the induction there is over the
derivation, and every clause of Definition~\ref{ne:def:elemplus} preserves dense
analyticity on its own.

\begin{theorem}\label{ne:thm:plusdense}
Every member of $\Elem^{-1}$ is analytic on a dense set.
\end{theorem}

\begin{proof}
Induction.  The elementary generators are \Cref{ne:thm:density}; the operations
are as in \Cref{ne:sec:density}; and the inverse rule splits on whether the test
set meets $U$.  Given a nonempty open $B$: if $B\cap U\ne\emptyset$ then
\Cref{ne:thm:invbranch} applies to the nonempty open set $B\cap U$ and produces a
point of $B$ at which $g$ is analytic; and if $B\cap U=\emptyset$ then
$B\subseteq U^{\mathsf c}$, so
$B\subseteq\operatorname{int}(U^{\mathsf c})$ because $B$ is open, and every
point of $B$ will do by the rule's analyticity clause.
\end{proof}

\begin{corollary}\label{ne:cor:plusmain}
Neither $F$ nor $F^{-1}$ agrees with any member of $\Elem^{-1}$ on any subset
of $[0,1]$ with nonempty interior.  In particular neither is representable on
$(0,1)$ by elementary functions, elementary operations, and any finite tower
of totalized inverse branches satisfying Definition~\ref{ne:def:elemplus}.
\end{corollary}

\begin{proof}
Let $h\in\Elem^{-1}$ and suppose it agrees with either $F$ or $F^{-1}$ on a
set $U\subseteq[0,1]$ with nonempty interior.  By
\Cref{ne:thm:plusdense}, $\An(h)$ is dense, hence meets
$\operatorname{int}U$.  At such a point the agreement gives equal germs, so
the corresponding Fabius function would be analytic there, contradicting
\Cref{ne:thm:nowhere} or \Cref{ne:thm:invnowhere}.
\end{proof}

\begin{boxedremark}
It is worth being explicit about what Corollary~\ref{ne:cor:plusmain} does and does
not rest on.  The inverse need not be two-sided: a continuous right inverse on
an open set is enough.  On the other hand, the conditional Lambert-$W$
constructor also requires the chosen totalization to be analytic on the
interior of that open set's complement.  The branch-domain boundary is
deliberately excluded from this side condition.  Because \texttt{Mathlib} does
not define $W$, the formalization neither constructs a standard real branch
nor verifies any particular totalization.
\end{boxedremark}
